\uuid{oStE}
\titre{DL usuels/Étude locale/Inflexion}
\auteur{hanani ; mzari}
\organisation{exo7}
\licence{CC BY-NC-SA 4.0}
\datecreate{2020-01-20}
\niveau{L1}
\module{Analyse}
\chapitre{Développement limité}
\sousChapitre{Applications}
\difficulte{2}
\theme{Applications des DL, DL usuels, Étude locale, Inflexion}
\exo7qcmid{lille-2.developpements-limites.034}

\contenu{
\texte{
Soit $a\in \Rr$ et $f$ une fonction. On dit que $f(x)$ est un $o\left((x-a)^n\right)$ au voisinage de $a$ et on écrit $f(x)=o((x-a)^n)$ si $\displaystyle \lim _{x\to a}\frac{f(x)}{(x-a)^n}=0$. On note par $DL_n(a)f(x)$ le développement limité de $f$ à l'ordre $n$ en $a$. On note aussi $DL_n(a^+)f(x)$ (resp. $DL_n(a^-)f(x)$) le développement limité de $f$ à l'ordre $n$ à droite en $a$ (resp. à gauche en $a$).
}

\begin{question}
\qclassification{125.03}
\qtags{motcle=DL usuels/Étude locale/Inflexion}
Soit $\displaystyle f(x)=\frac{x^2}{x^2+2x+2}$ et $g(x)=f(x-1)$. On note $T$ la tangente au graphe de $f$ au point d'abscisse $-1$. Parmi les affirmations suivantes, cocher celles qui sont vraies :
\begin{answers}
\good{$\displaystyle DL_3(0)g(x)=1-2x+2x^3+o(x^3)$.}
\bad{$T$ est la droite d'équation $y=1-2x$ et le graphe de $f$ est au dessus de $T$ au voisinage de $-1$.}
\bad{$T$ est la droite d'équation $y=1-2x$ et le graphe de $f$ est en dessous de $T$ au voisinage de $-1$.}
\bad{$T$ est la droite d'équation $y=1-2x$ et le point d'abscisse $-1$ est un point d'inflexion.}
\end{answers}
\begin{explanations}
D'abord $\displaystyle g(x)=\frac{1-2x+x^2}{1+x^2}=1-2x+2x^3+o(x^3)$. On en déduit que
$$\displaystyle DL_3(-1)f(x)=1-2(x+1)+2(x+1)^3+o\left((x+1)^3\right).$$
Donc $T$ est la droite d'équation $y=1-2(x+1)$ et le point d'abscisse $-1$ est un point d'inflexion.
\end{explanations}
\end{question}
}
